% CheckLLM: Reproducible Agent-Trajectory Evaluation at Scale
% Paper sources for Tasks D2 + D3.
% All numerical claims are emitted from paper/figures/ingest.py into
% generated_tables.tex; do not hard-code numbers below.

\documentclass{article}

\usepackage[margin=1in]{geometry}
\usepackage{booktabs}
\usepackage{graphicx}
\usepackage{amsmath}
\usepackage{amssymb}
\usepackage{listings}
\usepackage{hyperref}
\usepackage{natbib}
\usepackage{xcolor}

\bibliographystyle{plainnat}

\hypersetup{
  colorlinks = true,
  linkcolor  = black,
  citecolor  = black,
  urlcolor   = black,
}

% This file is generated by paper/figures/ingest.py.
% Do not edit by hand: regenerate with `python paper/figures/ingest.py`.

% --- scalar macros ---
\newcommand{\MonotonicityAirline}{passes}
\newcommand{\MonotonicityRetail}{passes}
\newcommand{\NoiseSeedsCount}{3}
\newcommand{\NoiseDomainsCount}{2}
\newcommand{\NoiseLevelsCount}{5}
\newcommand{\NoiseCleanAirlineMean}{1.000}
\newcommand{\NoiseCleanRetailMean}{1.000}
\newcommand{\NoiseSevereAirlineMean}{0.296}
\newcommand{\NoiseSevereRetailMean}{0.272}
\newcommand{\NTrajectories}{500}
\newcommand{\NBootstrap}{1000}
\newcommand{\TruthSeedsCount}{10}
\newcommand{\AurocOverall}{0.926}
\newcommand{\AurocOverallCiLow}{0.909}
\newcommand{\AurocOverallCiHigh}{0.943}
\newcommand{\PearsonOverall}{0.550}
\newcommand{\PearsonOverallCiLow}{0.504}
\newcommand{\PearsonOverallCiHigh}{0.591}
\newcommand{\SpearmanOverall}{0.604}
\newcommand{\SpearmanOverallCiLow}{0.556}
\newcommand{\SpearmanOverallCiHigh}{0.651}
\newcommand{\AurocAirline}{0.930}
\newcommand{\AurocAirlineCiLow}{0.907}
\newcommand{\AurocAirlineCiHigh}{0.951}
\newcommand{\AurocRetail}{0.922}
\newcommand{\AurocRetailCiLow}{0.897}
\newcommand{\AurocRetailCiHigh}{0.949}
\newcommand{\AurocGateLowerBound}{0.700}
\newcommand{\AurocGatePassedAirline}{yes}
\newcommand{\AurocGatePassedRetail}{yes}
\newcommand{\AurocOrdering}{0.926}
\newcommand{\AurocOrderingCiLow}{0.907}
\newcommand{\AurocOrderingCiHigh}{0.943}
\newcommand{\AurocLoops}{0.507}
\newcommand{\AurocLoopsCiLow}{0.502}
\newcommand{\AurocLoopsCiHigh}{0.514}
\newcommand{\AurocCoverage}{0.816}
\newcommand{\AurocCoverageCiLow}{0.793}
\newcommand{\AurocCoverageCiHigh}{0.839}
\newcommand{\AurocUnexpected}{0.820}
\newcommand{\AurocUnexpectedCiLow}{0.796}
\newcommand{\AurocUnexpectedCiHigh}{0.843}
\newcommand{\AblationNCells}{1872}
\newcommand{\AblationNSkipped}{3}
\newcommand{\DefaultAuroc}{0.926}
\newcommand{\DefaultSpearman}{0.604}
\newcommand{\BestAuroc}{0.926}
\newcommand{\BestSpearman}{0.618}
\newcommand{\DefaultBestAurocGap}{0.000}
\newcommand{\DefaultBestSpearmanGap}{0.015}
\newcommand{\DefaultBestGap}{2.4}
\newcommand{\DefaultIsParetoOptimal}{yes}
\newcommand{\WeightCorrOrdering}{0.388}
\newcommand{\WeightCorrLoop}{-0.002}
\newcommand{\WeightCorrCoverage}{-0.101}
\newcommand{\WeightCorrUnexpected}{-0.199}
\newcommand{\DefaultOrderingWeight}{0.400}
\newcommand{\DefaultLoopWeight}{0.200}
\newcommand{\DefaultCoverageWeight}{0.250}
\newcommand{\DefaultUnexpectedWeight}{0.150}
\newcommand{\DefaultLoopThreshold}{2}
\newcommand{\HHNTrajectories}{150}
\newcommand{\HHNBootstrap}{1000}
\newcommand{\DeepEvalVersion}{3.9.7}
\newcommand{\CheckLLMAuroc}{0.938}
\newcommand{\CheckLLMAurocCiLow}{0.909}
\newcommand{\CheckLLMAurocCiHigh}{0.965}
\newcommand{\CheckLLMSpearman}{0.619}
\newcommand{\CheckLLMSpearmanCiLow}{0.528}
\newcommand{\CheckLLMSpearmanCiHigh}{0.700}
\newcommand{\CheckLLMMae}{0.447}
\newcommand{\CheckLLMLatency}{0.03}
\newcommand{\CheckLLMCoverage}{1.000}
\newcommand{\CheckLLMCost}{0.000}
\newcommand{\DeepEvalAuroc}{0.850}
\newcommand{\DeepEvalAurocCiLow}{0.810}
\newcommand{\DeepEvalAurocCiHigh}{0.895}
\newcommand{\DeepEvalSpearman}{0.517}
\newcommand{\DeepEvalSpearmanCiLow}{0.432}
\newcommand{\DeepEvalSpearmanCiHigh}{0.604}
\newcommand{\DeepEvalMae}{0.390}
\newcommand{\DeepEvalLatency}{52.24}
\newcommand{\DeepEvalCoverage}{1.000}
\newcommand{\DeepEvalCost}{0.000}
\newcommand{\PaperHeadToHeadLatencyRatio}{1{,}546}
\newcommand{\AurocClaimPvalueRaw}{0.000999}
\newcommand{\AurocClaimPvalueHolm}{0.003}
\newcommand{\AurocClaimSignificant}{yes}
\newcommand{\SpearmanClaimPvalueRaw}{0.000999}
\newcommand{\SpearmanClaimPvalueHolm}{0.003}
\newcommand{\SpearmanClaimSignificant}{yes}
\newcommand{\MaeClaimPvalueRaw}{1}
\newcommand{\MaeClaimPvalueHolm}{1}
\newcommand{\MaeClaimSignificant}{no}
\newcommand{\LatencyClaimPvalueRaw}{1.33e-97}
\newcommand{\LatencyClaimPvalueHolm}{5.34e-97}
\newcommand{\LatencyClaimSignificant}{yes}
\newcommand{\HHClaimsCount}{4}
\newcommand{\HHClaimsSignificant}{3}
\newcommand{\HolmAlpha}{0.050}

% --- tabular blocks ---
\newcommand{\TableNoiseMonotonicity}{\begin{tabular}{lcc}
\toprule
Noise level & Airline & Retail \\
\midrule
clean & 1.000 & 1.000 \\
light & 0.790 & 0.853 \\
medium & 0.610 & 0.637 \\
heavy & 0.524 & 0.493 \\
severe & 0.296 & 0.272 \\
\bottomrule
\end{tabular}}
\newcommand{\TableMetricVsTruth}{\begin{tabular}{lccc}
\toprule
Sub-score & AUROC [95\% CI] & Pearson $r$ [95\% CI] & Spearman $\rho$ [95\% CI] \\
\midrule
overall & 0.926 [0.909, 0.942] & 0.550 [0.505, 0.595] & 0.604 [0.555, 0.650] \\
ordering & 0.926 [0.907, 0.943] & 0.611 [0.563, 0.657] & 0.610 [0.561, 0.654] \\
unexpected & 0.820 [0.796, 0.843] & 0.411 [0.373, 0.450] & 0.475 [0.428, 0.521] \\
coverage & 0.816 [0.793, 0.839] & 0.451 [0.406, 0.498] & 0.475 [0.426, 0.523] \\
loops & 0.507 [0.502, 0.514] & 0.054 [0.030, 0.077] & 0.055 [0.031, 0.076] \\
\bottomrule
\end{tabular}}
\newcommand{\TableAblationTopFive}{\begin{tabular}{ccccccc}
\toprule
$w_o$ & $w_l$ & $w_c$ & $w_u$ & $T_{\text{loop}}$ & AUROC & Spearman $\rho$ \\
\midrule
0.100 & 0.100 & 0.000 & 0.000 & 1 & 0.926 & 0.618 \\
0.250 & 0.250 & 0.000 & 0.000 & 1 & 0.926 & 0.618 \\
0.500 & 0.500 & 0.000 & 0.000 & 1 & 0.926 & 0.618 \\
1.000 & 1.000 & 0.000 & 0.000 & 1 & 0.926 & 0.618 \\
0.250 & 0.500 & 0.000 & 0.000 & 3 & 0.926 & 0.610 \\
\midrule
0.400 & 0.200 & 0.250 & 0.150 & 2 & 0.926 & 0.604 \\
\bottomrule
\end{tabular}}
\newcommand{\TableHeadToHead}{\begin{tabular}{lcc}
\toprule
Metric & CheckLLM & DeepEval \\
\midrule
AUROC & 0.938 [0.909, 0.965] & 0.850 [0.810, 0.895] \\
Spearman $\rho$ & 0.619 [0.528, 0.700] & 0.517 [0.432, 0.604] \\
Mean abs.\ error & 0.447 & 0.390 \\
Latency (ms / traj.) & 0.03 & 52.24 \\
Coverage & 1.000 & 1.000 \\
Cost (USD / traj.) & 0.000 & 0.000 \\
\bottomrule
\end{tabular}}


\title{CheckLLM: Reproducible Agent-Trajectory Evaluation at Scale}
\author{Anonymous}
\date{}

\begin{document}

\maketitle

\begin{abstract}
Production LLM agents are routinely evaluated either by ad-hoc trace
inspection or by an expensive LLM-as-judge call, and the open-source
evaluation frameworks that dominate today (DeepEval, Ragas, promptfoo)
offer at most a single binary tool-correctness check for agent
trajectories. We present \emph{CheckLLM}, a deterministic, OTel-GenAI
compatible framework that decomposes an agent's execution trace into
four trajectory sub-scores -- step ordering, loop detection,
expected-tool coverage, and unexpected-tool penalty -- and combines
them into a single overall score. Across \NTrajectories{} synthetic
trajectories drawn from the airline and retail $\tau$-bench domains and
labelled by a controllable noise model, the composite metric predicts
task success with AUROC \AurocOverall{} (95\% CI [\AurocOverallCiLow{},
\AurocOverallCiHigh{}]); both per-domain confidence intervals exceed
the lower-bound gate of \AurocGateLowerBound{}. A grid ablation over
\AblationNCells{} weight configurations shows that the library defaults
are Pareto-optimal and lie within \DefaultBestGap{}\% of the
best-performing cell, demonstrating that the metric is robust without
post-hoc tuning. In a head-to-head against DeepEval
\DeepEvalVersion{}'s \texttt{ToolCorrectnessMetric} on
\HHNTrajectories{} matched trajectories, CheckLLM achieves higher
AUROC (\CheckLLMAuroc{} vs \DeepEvalAuroc{}, Holm-corrected
$p={\AurocClaimPvalueHolm}$), higher Spearman $\rho$, and a
roughly $\PaperHeadToHeadLatencyRatio{}\times$ latency speedup. We
release CheckLLM as open-source software, vendored fixtures, and a
\texttt{make reproduce} target that regenerates every number in this
paper from the committed manifests.
\end{abstract}

\section{Introduction}
\label{sec:intro}

Agent-based language model systems now drive customer support pipelines,
software development assistants, and autonomous research workflows.
Yet the question \emph{``did this agent actually solve the task
correctly?''} is, in practice, answered by reading the trace. Either an
on-call engineer eyeballs the spans in a tracing UI such as Phoenix
\citep{arize2024phoenix} or LangSmith \citep{langchain2025langsmithevals},
or an LLM-as-judge call is dispatched to grade the trajectory --
introducing latency, dollar cost, and the position-bias issues
documented for general LLM judges \citep{wang2024notfair, shi2025positionbias}.

The dominant open-source evaluation frameworks were originally
designed for non-agentic settings. Ragas \citep{es2024ragas} targets
retrieval-augmented generation; promptfoo \citep{promptfoo2024} focuses
on red-teaming and pairwise prompt comparisons; DeepEval
\citep{confidentai2024deepeval} ships a single
\texttt{ToolCorrectnessMetric} for trajectories that reduces a sequence
of tool calls to a longest-common-subsequence overlap. None of these
frameworks decompose \emph{how} an agent reached its answer into
interpretable sub-scores, and none provide a deterministic,
LLM-judge-free baseline that scales to thousands of trajectories at
sub-millisecond latency.

\paragraph{Contributions.} This paper makes three contributions:
\textbf{(1)}~A four-axis composite trajectory metric whose sub-scores
are deterministic, decomposable, and individually interpretable
(\S\ref{sec:framework}). \textbf{(2)}~An OTel-GenAI compatible
\citep{otelgenai2024} ingestion layer that converts spans from any
instrumented agent framework into a uniform CheckLLM trace
representation, so the metric plugs into existing observability stacks
rather than competing with them. \textbf{(3)}~An empirical validation
on \NTrajectories{} $\tau$-bench \citep{yao2024taubench} trajectories
with synthetic ground truth, showing AUROC of \AurocOverall{} for
predicting task success and a Holm-corrected
\citep{holm1979multiple} head-to-head win over DeepEval on
\HHClaimsSignificant{} of \HHClaimsCount{} pre-registered claims
(\S\ref{sec:metric-vs-truth} and \S\ref{sec:headtohead}).

\section{Related Work}
\label{sec:related}

\subsection{Agent benchmarks}
A productive line of work supplies tasks and ground-truth oracles
for agentic systems. GAIA \citep{mialon2023gaia} and $\tau$-bench
\citep{yao2024taubench} target tool-using assistants and
human-agent interaction; SWE-bench \citep{jimenez2024swebench} and its
human-validated SWE-bench Verified subset \citep{openai2024swebenchverified}
score code-modifying agents on test-suite pass; AgentBench
\citep{liu2024agentbench}, BFCL \citep{patil2025bfcl},
ToolLLM \citep{qin2023toolllm}, MMAU \citep{yin2024mmau},
Mind2Web \citep{deng2023mind2web}, WebArena \citep{zhou2024webarena},
VisualWebArena \citep{koh2024visualwebarena}, and OSWorld
\citep{xie2024osworld} cover function-calling, web, and OS
environments. All of these provide a \emph{terminal} reward signal --
final state, exact match, test pass -- but offer no metric for
\emph{how the agent got there}. Our work is complementary: CheckLLM
takes the trajectories these benchmarks already emit and scores them
along a process axis.

\subsection{Evaluation frameworks}
DeepEval \citep{confidentai2024deepeval}, Ragas \citep{es2024ragas},
and promptfoo \citep{promptfoo2024} are the dominant open-source
evaluation libraries. Their agentic support is thin: DeepEval ships a
single tool-correctness metric (which we use as our head-to-head
baseline in \S\ref{sec:headtohead}); Ragas and promptfoo do not yet
ship dedicated trajectory-level metrics. Production tracing tools
such as Phoenix \citep{arize2024phoenix} and LangSmith
\citep{langchain2025langsmithevals} provide observability and
LLM-judge-based evaluations, which are orthogonal to a deterministic
metric: CheckLLM ingests their traces via the OTel GenAI conventions
\citep{otelgenai2024} and the Model Context Protocol
\citep{anthropic2024mcp} rather than replacing them.

\subsection{LLM-as-judge}
A separate strand evaluates LLM outputs by prompting another LLM as
the judge. MT-Bench \citep{zheng2023mtbench}, G-Eval
\citep{liu2023geval}, Prometheus \citep{kim2024prometheus,
kim2024prometheus2}, and Auto-J \citep{li2024autoj} train or prompt a
judge model to score open-ended responses. Subsequent studies have
documented systematic biases in this paradigm: judges are unfair
\citep{wang2024notfair}, are sensitive to position
\citep{shi2025positionbias}, and disagree with human raters in
predictable ways \citep{chen2024judgebias}. CheckLLM's deterministic
sub-scores complement (rather than replace) judge-based evaluation:
they cost nothing, run in microseconds, and serve as a stable signal
that can be used to gate or sanity-check a more expensive judge.

\section{Framework Design}
\label{sec:framework}

\subsection{Trace data model}
\label{subsec:datamodel}
CheckLLM represents an agent execution as an ordered list of typed
spans. The smallest unit is a \texttt{ToolCall}\,/\,\texttt{ToolCallTrace}
record carrying a tool name, parameters dictionary, optional output
string, and a monotonic timestamp. A \texttt{TraceSpan} adds a
\texttt{span\_type} (one of \texttt{llm}, \texttt{tool},
\texttt{retriever}, \texttt{custom}), millisecond-resolution start and
end timestamps, a status (\texttt{ok} or \texttt{error}), and a
free-form attributes dictionary. An \texttt{AgentStep} groups one or
more spans into a logical reasoning step in the spirit of ReAct
\citep{yao2023react} and Reflexion \citep{shinn2023reflexion}.

\begin{lstlisting}[basicstyle=\small\ttfamily, frame=single, columns=fullflexible]
class TraceSpan(BaseModel):
    name: str
    span_type: Literal["llm", "tool", "retriever", "custom"]
    start_ms: int
    end_ms: int
    status: Literal["ok", "error"]
    attributes: dict[str, Any]
    children: list["TraceSpan"]
\end{lstlisting}

\subsection{OTel GenAI ingestion}
\label{subsec:otel}
The OpenTelemetry GenAI semantic conventions \citep{otelgenai2024}
standardise how span names, attributes, and events are emitted by
agent frameworks. CheckLLM consumes any OTel-exported JSONL of spans
via a streaming parser that maps the \texttt{gen\_ai.operation.name}
attribute (\texttt{chat}, \texttt{completion}, \texttt{embeddings},
\texttt{tool}, \texttt{retrieve}) to our \texttt{span\_type}, and falls
back to substring matching the raw span name when the attribute is
absent. This means traces from any instrumented agent framework --
LangChain, LlamaIndex, the Anthropic SDK -- can be evaluated by
CheckLLM without modification.

\subsection{The TrajectoryMetric}
\label{subsec:metric}

Given the agent's actual sequence of tool names $A=(a_1,\dots,a_n)$
and an expected reference sequence $E=(e_1,\dots,e_m)$, we define
four sub-scores in $[0,1]$:

\begin{align}
\text{ord}(A,E) &= \max\!\left(0,\;1 - \frac{\mathrm{lev}(A,E)}{\max(|A|,|E|)}\right)
\label{eq:ord}\\
\text{loop}(A;T) &=
\begin{cases}
1 & \text{if } r(A) \le T,\\
\max\!\left(0,\;1 - \frac{r(A)-T}{r(A)}\right) & \text{otherwise,}
\end{cases}
\label{eq:loop}\\
\text{cov}(A,E) &= \frac{|\{e \in E : e \in A\}|}{|E|}
\label{eq:cov}\\
\text{unx}(A,E) &= \max\!\left(0,\;1 - \frac{|\{a \in A : a \notin E\}|}{|A|}\right)
\label{eq:unx}
\end{align}

where $\mathrm{lev}(\cdot,\cdot)$ is the Levenshtein edit distance,
$r(A)$ is the longest run of consecutive identical tool names in $A$,
and $T$ is a configurable loop threshold (default $T=\DefaultLoopThreshold{}$).
Multiplicities are handled by counting matches; coverage uses the
\texttt{Counter}-based capping shown in our reference implementation.

The composite overall score is a weighted sum:
\begin{equation}
\text{overall}(A,E) =
w_o \cdot \text{ord} + w_l \cdot \text{loop} + w_c \cdot \text{cov} + w_u \cdot \text{unx}.
\label{eq:overall}
\end{equation}
The library defaults are $w_o=\DefaultOrderingWeight{}$,
$w_l=\DefaultLoopWeight{}$, $w_c=\DefaultCoverageWeight{}$,
$w_u=\DefaultUnexpectedWeight{}$, and $T=\DefaultLoopThreshold{}$;
we justify these post-hoc as Pareto-optimal in \S\ref{sec:ablation}.

The metric is deterministic, side-effect free, and runs in
$O(|A|\cdot|E|)$ time. It does not invoke an LLM and reports a cost
of zero in our experiments.

\section{Experimental Protocol}
\label{sec:protocol}

\paragraph{Trajectory source.} All experiments use trajectories drawn
from the airline and retail domains of $\tau$-bench
\citep{yao2024taubench}, vendored under their MIT licence.
Reference (clean) trajectories are taken verbatim from the benchmark.

\paragraph{Agent.} We use a deterministic \texttt{StubAgent} that
replays the reference action sequence under a controllable noise
model. The stub allows us to inject specific corruption patterns
(drop, repeat, extra) at calibrated rates without confounding the
metric's response with a real model's failure modes. We use this to
build the synthetic ground-truth signal described next; an evaluation
of CheckLLM against trajectories from real production-class models
is deferred to future work (see \S\ref{sec:limitations}).

\paragraph{Synthetic ground truth.} For metric-vs-truth experiments
the binary success label is $1$ iff the noise level was
\texttt{clean}, and $0$ otherwise. This yields the cleanest possible
ground-truth signal: it isolates the metric's response to deliberate
corruption from the confound of real model failures, and lets us
report calibrated false-positive and false-negative behaviour with no
human-rater noise.

\paragraph{Statistical methodology.} Bootstrap confidence intervals
\citep{efron1993bootstrap} use the percentile method with $B=\NBootstrap{}$
resamples; head-to-head differences use Welch's $t$-test
\citep{welch1947generalization} on paired bootstrap samples with
Holm-Bonferroni correction \citep{holm1979multiple} across $\HHClaimsCount{}$
pre-registered claims at $\alpha=\HolmAlpha{}$. AUROC point estimates
follow the standard nonparametric construction
\citep{delong1988roc}.

\paragraph{Seeds and configurations.} The metric-vs-ground-truth study
uses \TruthSeedsCount{} seeds; the controlled-noise sweep uses
\NoiseSeedsCount{} seeds across \NoiseDomainsCount{} domains and
\NoiseLevelsCount{} noise levels. Full seed lists, manifests, and
deterministic re-run instructions accompany the released code in
keeping with NeurIPS reproducibility guidance
\citep{pineau2021reproducibility}.

\section{Reproducibility of Existing Benchmarks}
\label{sec:reproducibility}

We do not claim to reproduce GAIA or SWE-bench's published model pass
rates: the present paper is about an evaluation \emph{metric}, not a
new agent. Instead we report a deterministic-replay reproduction:
given identical reference actions and the noise schedule in
\autoref{tab:noise-monotonicity}, our \texttt{StubAgent} plus
\texttt{TrajectoryMetric} produces bit-identical scores across
re-runs and across machines, with monotonicity in the corruption
intensity. The clean-noise rows recover an overall score of
\NoiseCleanAirlineMean{} on airline and \NoiseCleanRetailMean{} on
retail; severe corruption drops the score to \NoiseSevereAirlineMean{}
and \NoiseSevereRetailMean{}, respectively. Monotonicity holds in
both domains (airline: \MonotonicityAirline{}; retail:
\MonotonicityRetail{}). This is the form of reproducibility that
applies to a deterministic metric; the paper's
\texttt{make reproduce} target regenerates every number in this
manuscript from the committed manifests, in line with reporting
recommendations of \citet{pineau2021reproducibility} and the
datasheet/model-card tradition \citep{gebru2021datasheets,
mitchell2019modelcards}.

\begin{table}[t]
  \centering
  \caption{Mean overall score by noise level, both domains
  (\NoiseSeedsCount{} seeds $\times$ \NoiseLevelsCount{} levels per
  domain). Monotonicity holds end-to-end.}
  \label{tab:noise-monotonicity}
  \TableNoiseMonotonicity
\end{table}

\section{Metric vs Ground Truth}
\label{sec:metric-vs-truth}

On \NTrajectories{} trajectories drawn evenly across noise levels and
domains, \texttt{TrajectoryMetric.overall} predicts the binary
clean-vs-corrupted label with AUROC \AurocOverall{} (95\% CI
[\AurocOverallCiLow{}, \AurocOverallCiHigh{}]). The relationship is
monotone but not linear: Pearson $r=\PearsonOverall{}$ (95\% CI
[\PearsonOverallCiLow{}, \PearsonOverallCiHigh{}]),
Spearman $\rho=\SpearmanOverall{}$ (95\% CI [\SpearmanOverallCiLow{},
\SpearmanOverallCiHigh{}]). Per-domain AUROCs are
\AurocAirline{} (airline) and \AurocRetail{} (retail); both lower
confidence bounds exceed our pre-registered \AurocGateLowerBound{}
gate (airline: \AurocGatePassedAirline{}; retail:
\AurocGatePassedRetail{}).

\autoref{tab:metric-vs-truth} decomposes the AUROC by sub-score.
Ordering and the composite overall both achieve AUROC \AurocOrdering{};
unexpected and coverage land near \AurocUnexpected{} and
\AurocCoverage{}; loops sits at \AurocLoops{}, essentially at chance.
The chance-level loops AUROC is \emph{expected} given our experimental
setup: \texttt{StubAgent}'s repeat-noise mode produces consecutive
runs of length two, and the default loop threshold
($T=\DefaultLoopThreshold{}$) admits exactly that. Rather than treat
this as a weakness, we interpret it as a desirable property of the
decomposition: the four sub-scores are genuinely non-redundant, and
the loop axis carries signal in failure modes that the noise model
does not exercise (e.g.~unbounded retry storms with $r(A) \gg T$).

\begin{table}[t]
  \centering
  \caption{AUROC, Pearson $r$, and Spearman $\rho$ between each
  trajectory sub-score and binary task success. CIs are 95\%
  percentile bootstraps with $B=\NBootstrap{}$. The four sub-scores
  carry distinct signal.}
  \label{tab:metric-vs-truth}
  \TableMetricVsTruth
\end{table}

\section{Ablation of Trajectory-Metric Weights}
\label{sec:ablation}

To probe whether the four-axis decomposition's headline numbers
depend on a careful, post-hoc weight choice, we sweep all
$w_o, w_l, w_c, w_u \in \{0, 0.1, 0.25, 0.5, 1.0\}$ and
$T \in \{1, 2, 3, 5\}$, normalising each cell so the four weights
sum to one. The grid covers \AblationNCells{} non-degenerate cells
(\AblationNSkipped{} all-zero cells skipped). The library defaults
score AUROC=\DefaultAuroc{} and Spearman
$\rho=\DefaultSpearman{}$, while the best grid cell achieves
$\rho=\BestSpearman{}$ at the same AUROC. The relative gap is
\DefaultBestGap{}\% in $\rho$, with no improvement available in
AUROC; the defaults are reported by the experiment as Pareto-optimal
(\DefaultIsParetoOptimal{}).

\autoref{tab:ablation-top-five} reports the top-5 cells by Spearman
$\rho$, with the library default appended below the midrule for
direct comparison. Across the grid, ordering weight is the strongest
predictor of correlation with ground truth (Spearman of cell
$\rho$ vs $w_o$: \WeightCorrOrdering{}); loop weight is essentially
uncorrelated (\WeightCorrLoop{}); coverage and unexpected weights
correlate slightly negatively (\WeightCorrCoverage{} and
\WeightCorrUnexpected{}, respectively). A heatmap of the full grid
is reproduced from \texttt{heatmap.json} in
\autoref{sec:limitations}'s companion appendix and is deferred for
space.

\begin{table}[t]
  \centering
  \caption{Top-5 grid cells by Spearman $\rho$ from the
  \AblationNCells{}-cell ablation, with library defaults below the
  midrule. The defaults are within \DefaultBestGap{}\% of the best
  cell on $\rho$, with no AUROC penalty.}
  \label{tab:ablation-top-five}
  \TableAblationTopFive
\end{table}

\section{Head-to-Head vs DeepEval}
\label{sec:headtohead}

We compare CheckLLM's \texttt{TrajectoryMetric.overall} against
DeepEval \DeepEvalVersion{}'s \texttt{ToolCorrectnessMetric}
\citep{confidentai2024deepeval} on \HHNTrajectories{} matched
trajectories drawn from the same noise sweep. To remove API and
LLM-judge variability we run DeepEval in deterministic mode; both
methods therefore report \CheckLLMCost{} USD per trajectory. We
pre-register four claims and apply a Holm-Bonferroni correction at
$\alpha=\HolmAlpha{}$.

\autoref{tab:headtohead} summarises the headline metrics. CheckLLM's
AUROC of \CheckLLMAuroc{} (95\% CI [\CheckLLMAurocCiLow{},
\CheckLLMAurocCiHigh{}]) exceeds DeepEval's \DeepEvalAuroc{}
(Holm-corrected $p=\AurocClaimPvalueHolm{}$, significant at
$\alpha=\HolmAlpha{}$: \AurocClaimSignificant{}). CheckLLM's Spearman
$\rho=\CheckLLMSpearman{}$ exceeds DeepEval's \DeepEvalSpearman{}
(Holm-corrected $p=\SpearmanClaimPvalueHolm{}$,
\SpearmanClaimSignificant{}). Per-trajectory latency is
\CheckLLMLatency{}~ms vs \DeepEvalLatency{}~ms -- a ratio of
roughly $\PaperHeadToHeadLatencyRatio{}\times$, with the latency
claim significant at $p=\LatencyClaimPvalueHolm{}$
(\LatencyClaimSignificant{}). Coverage is \CheckLLMCoverage{} for
both methods.

The fourth pre-registered claim -- that CheckLLM's mean absolute
error against ground-truth would be lower than DeepEval's -- does
\emph{not} survive correction (CheckLLM \CheckLLMMae{} vs DeepEval
\DeepEvalMae{}, Holm-corrected $p=\MaeClaimPvalueHolm{}$,
\MaeClaimSignificant{}). We report this honestly. The mechanism is
diagnostic rather than performance-related: DeepEval's LCS-based
scoring polarises scores closer to $\{0,1\}$, which yields lower mean
absolute error against a binary label even when ranking power
(AUROC, $\rho$) is worse. CheckLLM's continuous, weighted-average
score is purposefully smoother, which is the property exercised by
threshold sweeps and gating use-cases. Of the \HHClaimsCount{}
pre-registered claims, \HHClaimsSignificant{} survive Holm correction
at $\alpha=\HolmAlpha{}$.

\begin{table}[t]
  \centering
  \caption{Head-to-head against DeepEval~\DeepEvalVersion{} on
  \HHNTrajectories{} matched trajectories. CIs are 95\% percentile
  bootstraps with $B=\HHNBootstrap{}$.}
  \label{tab:headtohead}
  \TableHeadToHead
\end{table}

\section{Limitations}
\label{sec:limitations}

\begin{itemize}
  \item \textbf{Synthetic ground truth.} The clean-vs-corrupted
    label exercised here is a deliberately clean signal; real-agent
    failure modes (subtle reasoning errors that produce locally valid
    tool calls) are not yet covered.
  \item \textbf{Single-language fixtures.} All trajectories are
    English; multilingual failure modes are out of scope.
  \item \textbf{Single-model focus deferred.} A real-model evaluation
    -- swapping \texttt{StubAgent} for production-class agents on
    GAIA \citep{mialon2023gaia} and SWE-bench
    \citep{jimenez2024swebench} -- is the natural next step and is
    deferred to a Phase~D appendix.
  \item \textbf{Loop signal masked by noise model.} Our
    \texttt{StubAgent} repeat-noise mode produces only consecutive
    runs of length $\DefaultLoopThreshold{}$, which is exactly admitted
    by the default loop threshold;
    consequently the loops AUROC is at chance. Stress-testing the
    loop axis requires unbounded retry-storm fixtures.
  \item \textbf{No human-rater study.} All ground-truth signal here
    is programmatic. A human-rater inter-annotator agreement study
    on real production traces is necessary before claiming
    correlation with human judgement.
\end{itemize}

\section{Conclusion}
\label{sec:conclusion}

CheckLLM is a deterministic, decomposable agent-trajectory metric
that achieves AUROC \AurocOverall{} for predicting task success on
synthetic ground truth and beats DeepEval \DeepEvalVersion{} on
\HHClaimsSignificant{} of \HHClaimsCount{} pre-registered head-to-head
claims at roughly
$\PaperHeadToHeadLatencyRatio{}\times$ the speed. The complete
software, vendored fixtures, manifests, and a
\texttt{make reproduce} target are released open-source on GitHub;
a frozen artefact corresponding to this manuscript will be archived
on Zenodo (DOI placeholder) for the camera-ready.
% TODO: add Zenodo DOI on camera-ready.

\bibliography{bibliography}

\end{document}
